\metatitle{P-Grothendieck polynomials}
\metadescription{An introduction to P-Grothendieck polynomials, their definition via shifted set-valued tableaux, and their relation to weak P-Grothendieck polynomials and Schur-P positivity.}



\section[pGrothendieck]{P-Grothendieck polynomials}

\begin{polydata}{pGrothendieck}
  Name   & P-Grothendieck polynomials \\
  Space  & Sym \\
  Basis  & True \\
  Rating & 1 \\
  Bib    & Hawkes2023 \\
  Year   & 2023 \\
  Category & Schur \\
  KTheoreticAnalogueOf & schurP | Hawkes2023 \\
  PositiveIn & schurP | Hawkes2023 | scope=weak P-Grothendieck polynomials \\
\end{polydata}

The $P$-Grothendieck polynomials are symmetric, non-homogeneous symmetric
functions attached to strict partitions.  In \cite{Hawkes2023},
\name{Graham Hawkes} introduces weak $P$-Grothendieck polynomials, which play
the same role for \hyperref[schurP]{Schur $P$ functions} as
\hyperref[grothendieckDualStable]{dual stable Grothendieck polynomials} play
for Schur functions.  The combinatorial model uses \defin{shifted multiset
tableaux}: shifted tableaux whose boxes contain nonempty multisets in the
alphabet $1'\lt 1\lt 2'\lt 2\lt \dotsb$, with primed entries appearing with
multiplicity at most one in each box and with shifted row and column
semistandard conditions.

Following \cite[Def.~4.5]{Hawkes2023}, write $a\lt_u z$ if $a\lt z$, or if
$a=z$ and both entries are unprimed.  Similarly, write $a\lt_p z$ if
$a\lt z$, or if $a=z$ and both entries are primed.  If a box $b'$ is directly
to the right of a box $b$, then every entry $a\in b$ must satisfy
$a\lt_u z$ for every $z\in b'$.  If $b'$ is directly below $b$, then for
every $z\in b'$ there must exist some $a\in b$ with $a\lt_p z$.  The set
$\mathrm{SMT}_0(\mu)$ consists of those shifted multiset tableaux for which
the smallest entry in each row is unprimed.

Hawkes proves that these functions are \hyperref[schurP]{Schur $P$}-positive
by homogeneous degree.  This is analogous to the Schur-positivity of weak
stable Grothendieck polynomials.
The combinatorial definition is
\[
  \grothendieckPDual_\mu(\xvec;\tvec)
  =
  \sum_{P\in \mathrm{SMT}_0(\mu)}
  \tvec^{\mathrm{dw}(P)}\xvec^{\mathrm{wt}(P)},
\]
where $\mathrm{wt}(P)$ counts the values in $P$, ignoring primes, and
$\mathrm{dw}(P)$ records how many extra entries lie on each shifted diagonal;
see \cite[Thm.~4.17]{Hawkes2023}.

\begin{example*}[Shifted multiset tableaux for $\mu=21$]
% Related Rust: sym-poly/sym/src/p_grothendieck.rs and
% sym-poly/sym/examples/p_grothendieck_site_example.rs
Consider Hawkes's $t$-deformation for $\mu=21$ in two variables
$x_1,x_2$.  The first three homogeneous pieces are
\begin{align*}
\left[\deg_{\xvec}=3\right]\grothendieckPDual_{21}
&= x_1^2x_2+x_1x_2^2, \\
\left[\deg_{\xvec}=4\right]\grothendieckPDual_{21}
&= (t_1+t_2)
\left(x_1^3x_2+2x_1^2x_2^2+x_1x_2^3\right), \\
\left[\deg_{\xvec}=5\right]\grothendieckPDual_{21}
&= (t_1^2+t_2^2)
\left(x_1^4x_2+2x_1^3x_2^2+2x_1^2x_2^3+x_1x_2^4\right) \\
&\quad
+ t_1t_2
\left(x_1^4x_2+3x_1^3x_2^2+3x_1^2x_2^3+x_1x_2^4\right).
\end{align*}
The degree-four part comes from the eight tableaux in $\mathrm{SMT}_0(21)$
with exactly one extra entry beyond the three boxes of the shifted shape.
Equivalently, this degree-four piece is
$(t_1+t_2)\schurP_{31}(x_1,x_2)$; compare \cite[Ex.~4.18]{Hawkes2023}.
\end{example*}

Suppose $\mu$ is a strict partition with $m$ parts.
The \defin{weak symmetric $P$-Grothendieck polynomial} $\grothendieckPDual_\mu(\xvec)$
in $n \geq m$ variables is defined as
\[
\grothendieckPDual_\mu(\xvec) =
\prod_{i\lt j} (x_i-x_j)^{-1}
\sum_{\sigma \in \symS_n / \symS_{n-m}}
\sign(\sigma)
\left(
\prod_{i=1}^m
\left(
\frac{x_{\sigma_i}}{1-x_{\sigma_i}}
\right)^{\mu_i}
\right)
\left(
\prod_{i \lt j, i \leq m}
x_{\sigma_i} + x_{\sigma_j}
\right)
\left(
\prod_{m \lt i \lt j}
x_{\sigma_i} - x_{\sigma_j}
\right).
\]
Here, $\symS_n / \symS_{n-m}$ is the set of permutations in $\symS_n$
with no descent after position $m$.
